\documentclass[a4paper]{article}
\usepackage{fancyhdr}
\usepackage{tocloft}
\usepackage{titlesec}
\usepackage{listings}
\usepackage{color}
\usepackage{graphicx} % For including graphics (images)
\usepackage{amsfonts} % For using additional math fonts
\usepackage[table]{xcolor} % For coloring tables
\usepackage{pagecolor} % For setting the background color of the page
\usepackage{comment} % For including multi-line comments
\usepackage{amsmath} % For advanced mathematical typesetting
\usepackage{tcolorbox} % For creating colored and framed text boxes
\usepackage{tikz} % For creating graphics programmatically
\usepackage{multicol} % For creating multi-column layouts
\usepackage[british]{babel} % For setting the document language to British English
\usepackage[left = 2cm, right = 2cm, top = 3cm, bottom = 3cm]{geometry} % For setting page margins
\usepackage{marvosym} % For providing icons for contact details
\usepackage{enumitem} % For redefining spacing in lists
\usepackage{FiraSans} % For changing the document font to Fira Sans
\usepackage{pgfplots} % For creating plots and graphs
\usepackage{paracol} % For typesetting columns of text in parallel
\usepackage{calc} % For performing calculations in LaTeX
\usepackage{pgffor} % For using loops in LaTeX
\usepackage{ifthen} % Provides the \ifthenelse command, which allows you to evaluate logical expressions and execute different code paths based on whether the expression is true or false.
\usepackage{xparse}
\usepackage{expl3}
\usepackage{bchart}
\usepackage{ulem}
\usepackage[colorlinks = true, urlcolor = white, linkcolor = white]{hyperref} % For creating hyperlinks with colored text

% Dark Mode Color Pallet
\definecolor{palette.dark.back}{HTML}{1A1A1D}
\definecolor{palette.dark.frame}{HTML}{49494C}
\definecolor{palette.dark.title}{HTML}{3b3b3c}
\definecolor{palette.dark.text}{HTML}{FFFFFF}

% Light Mode Color Pallet
\definecolor{palette.light.back}{HTML}{F0F0F0}
\definecolor{palette.light.frame}{HTML}{D0D0D0}
\definecolor{palette.light.title}{HTML}{5A5A5A}
\definecolor{palette.light.text}{HTML}{2A2A2A}

% Success Color Pallet
\definecolor{palette.success.text}{HTML}{DFF7E0}
\definecolor{palette.success.frame}{HTML}{81c784}
\definecolor{palette.success.title}{HTML}{66bb6a}
\definecolor{palette.success.back}{HTML}{388e3c}

% Info Color Pallet
\definecolor{palette.info.frame}{HTML}{4fc3f7}
\definecolor{palette.info.title}{HTML}{29b6f6}
\definecolor{palette.info.back}{HTML}{0288d1}
\definecolor{palette.info.text}{HTML}{B3E5FC}

% Warning Color Pallet
\definecolor{palette.warning.frame}{HTML}{ffb74d}
\definecolor{palette.warning.title}{HTML}{ffa726}
\definecolor{palette.warning.back}{HTML}{f57c00}
\definecolor{palette.warning.text}{HTML}{FFF3E0}

% Error Color Pallet
\definecolor{palette.error.frame}{HTML}{e57373}
\definecolor{palette.error.title}{HTML}{f44336}
\definecolor{palette.error.back}{HTML}{d32f2f}
\definecolor{palette.error.text}{HTML}{FFEBEE}

% Primary Color Pallet
\definecolor{palette.primary.frame}{HTML}{e3f2fd}
\definecolor{palette.primary.title}{HTML}{90caf9}
\definecolor{palette.primary.back}{HTML}{42a5f5}
\definecolor{palette.primary.text}{HTML}{BBDEFB}

% Secondary Color Pallet
\definecolor{palette.secondary.frame}{HTML}{f3e5f5}
\definecolor{palette.secondary.title}{HTML}{ce93d8}
\definecolor{palette.secondary.back}{HTML}{ab47bc}
\definecolor{palette.secondary.text}{HTML}{EDE7F6}

\fancyhf{}
\fancyhead[L]{POSIX-Nexus Documentation}
\fancyhead[R]{\thepage}
\fancyfoot[C]{\textbf{POSIX-Nexus}}
\renewcommand{\headrulewidth}{0.4pt}
\renewcommand{\footrulewidth}{0.4pt}
\pagestyle{fancy}
\titleformat{\section}[block]{\Large\bfseries}{\thesection}{1em}{}
\titleformat{\subsection}[block]{\large\bfseries}{\thesubsection}{1em}{}
\titleformat{\subsubsection}[block]{\normalsize\bfseries}{\thesubsubsection}{1em}{}
\title{
    \includegraphics[width=0.5\textwidth]{images/posix-nexus-icon.png}\\[1cm]
    \textbf{POSIX-Nexus: Enhancing POSIX Shell Performance}
}

\author{Canine-Table}
\date{\today}
\begin{document}
\maketitle
\begin{abstract}
POSIX-Nexus aims to enhance the performance of the POSIX shell by leveraging the processing capabilities of POSIX \texttt{awk} and \texttt{sed} backends. This document provides an overview, usage instructions, and contribution guidelines.
\end{abstract}
\newpage
\tableofcontents
\newpage
\newpage
\section{Introduction}
\label{Introduction}
\begin{NexMainBox}[light, hdrA, sdwA, crnA, grwB, secA]
	\begin{NexMainBox}[dark, crnA]
	\end{NexMainBox}
	\begin{NexMainBox}[dark, crnA]
		\begin{NexListDark}
			\NexItemDark{\NexLink{Network Types}{Network Types}: Categorizes networks based on their scope, size, and purpose. Common types include LAN, WAN, MAN, and PAN, each with unique characteristics and use cases.}
			\NexItemDark{\NexLink{Broadband Technologies}{Broadband Technologies}: Encompasses various technologies such as DSL, Cable, Fiber-optic, Satellite, and Mobile broadband, each offering unique advantages based on speed, availability, and cost.}
		\end{NexListDark}
	\end{NexMainBox}
\end{NexMainBox}

\newpage
\subsection{Network Types}
\label{Network Types}
\begin{NexMainBox}[light, hdrA, sdwA, crnA, grwB, ssecA]
	\begin{NexMainBox}[dark, crnA]
		Networks are broadly categorized based on their geographical scope, size, and purpose. This categorization helps determine the appropriate technologies, devices, and configurations for specific networking needs.
	\end{NexMainBox}
	\begin{NexMainBox}[dark, crnA]
		\begin{NexListDark}
			\NexItemDark{\NexLink{PAN (Personal Area Network)}{PAN (Personal Area Network)}: Designed for individual use within a limited range, including Bluetooth or USB connections.}
			\NexItemDark{\NexLink{LAN (Local Area Network)}{LAN (Local Area Network)}: Covers small, localized areas such as homes, schools, or offices. Typically relies on Ethernet or Wi-Fi.}
			\NexItemDark{\NexLink{MAN (Metropolitan Area Network)}{MAN (Metropolitan Area Network)}: Covers regions like cities or campuses, bridging multiple LANs within a metropolitan area.}
			\NexItemDark{\NexLink{WAN (Wide Area Network)}{WAN (Wide Area Network)}: Spans large geographical areas and connects multiple LANs. The internet is the largest example.}
			\NexItemDark{\NexLink{GAN (Global Area Network)}{GAN (Global Area Network)}: A type of network that spans across the globe, connecting multiple WANs and providing international connectivity. Commonly used by multinational organizations and global services.}
		\end{NexListDark}
	\end{NexMainBox}
\end{NexMainBox}

\newpage
\subsubsection{PAN (Personal Area Network)}
\label{PAN (Personal Area Network)}
\begin{NexMainBox}[light, hdrA, sdwA, crnA, grwB, sssecA]
	\begin{NexMainBox}[dark, crnA]
		A Personal Area Network (PAN) is a short-range network designed for individual use, often connecting personal devices such as smartphones, laptops, and wearable devices.
	\end{NexMainBox}
	\begin{NexMainBox}[dark, crnA]
		\begin{NexListDark}
			\NexItemDark{\NexOption{Scope}: Limited to an individual's workspace, typically within a 10-meter range.}
			\NexItemDark{\NexOption{Applications}:
				\begin{NexListLight}
					\NexItemLight{Facilitates Bluetooth connections for device pairing.}
					\NexItemLight{Used for connecting peripherals like keyboards, mice, and printers.}
					\NexItemLight{Supports data transfer between personal devices.}
				\end{NexListLight}
			}
			\NexItemDark{\NexOption{Technology}: Commonly uses Bluetooth, NFC, or USB for connectivity.}
		\end{NexListDark}
	\end{NexMainBox}
\end{NexMainBox}

\newpage
\subsubsection{LAN (Local Area Network)}
\label{LAN (Local Area Network)}
\begin{NexMainBox}[light, hdrA, sdwA, crnA, grwB, sssecA]
		\begin{NexMainBox}[dark, crnA]
			A Local Area Network (LAN) is a network that spans a small geographical area, such as a home, office, or school. It is ideal for connecting devices within close proximity, typically using Ethernet or Wi-Fi.
		\end{NexMainBox}
		\begin{NexMainBox}[dark, crnA]
			\begin{NexListDark}
			\NexItemDark{\NexOption{Scope}: Covers limited areas such as buildings or small campuses.}
			\NexItemDark{\NexOption{Applications}:
				\begin{NexListLight}
					\NexItemLight{Used in offices for sharing printers, files, and internet access.}
					\NexItemLight{Commonly set up in homes for device connectivity.}
					\NexItemLight{Facilitates gaming networks for multiplayer gaming.}
				\end{NexListLight}
			}
			\NexItemDark{\NexOption{Technology}: Implements Ethernet (wired) or Wi-Fi (wireless) protocols for connectivity.}
		\end{NexListDark}
	\end{NexMainBox}
\end{NexMainBox}

\newpage
\subsubsection{MAN (Metropolitan Area Network)}
\label{MAN (Metropolitan Area Network)}
\begin{NexMainBox}[light, hdrA, sdwA, crnA, grwB, sssecA]
	\begin{NexMainBox}[dark, crnA]
		A Metropolitan Area Network (MAN) connects multiple LANs within a metropolitan area, such as a city or a large campus. MANs bridge the gap between LANs and WANs for broader regional connectivity.
	\end{NexMainBox}
	\begin{NexMainBox}[dark, crnA]
		\begin{NexListDark}
			\NexItemDark{\NexOption{Scope}: Covers metropolitan regions, such as cities or extensive campuses.}
			\NexItemDark{\NexOption{Applications}:
				\begin{NexListLight}
					\NexItemLight{Used by city governments for public services like surveillance and internet access.}
					\NexItemLight{Facilitates inter-campus connectivity for universities.}
					\NexItemLight{Enables metro Ethernet services for business districts.}
				\end{NexListLight}
			}
			\NexItemDark{\NexOption{Technology}: Employs fiber optic networks and metro Ethernet for high-speed connectivity.}
		\end{NexListDark}
	\end{NexMainBox}
\end{NexMainBox}

\newpage
\subsubsection{WAN (Wide Area Network)}
\label{WAN (Wide Area Network)}
\begin{NexMainBox}[light, hdrA, sdwA, crnA, grwB, sssecA]
	\begin{NexMainBox}[dark, crnA]
		A Wide Area Network (WAN) spans large geographical areas, connecting multiple LANs. WANs are essential for communication between regions and are the foundation for the global internet.
	\end{NexMainBox}
	\begin{NexMainBox}[dark, crnA]
		\begin{NexListDark}
			\NexItemDark{\NexOption{Scope}: Extends across cities, countries, or continents.}
			\NexItemDark{\NexOption{Applications}:
				\begin{NexListLight}
					\NexItemLight{Used by enterprises to connect offices across regions.}
					\NexItemLight{Supports internet service providers (ISPs) to deliver connectivity.}
					\NexItemLight{Provides connectivity for large-scale communication systems.}
				\end{NexListLight}
			}
			\NexItemDark{\NexOption{Technology}: Utilizes fiber optic cables, satellite links, and dedicated leased lines.}
		\end{NexListDark}
	\end{NexMainBox}
\end{NexMainBox}

\newpage
\subsubsection{GAN (Global Area Network)}
\label{GAN (Global Area Network)}
\begin{NexMainBox}[light, hdrA, sdwA, crnA, grwB, sssecA]
	\begin{NexMainBox}[dark, crnA]
	A Global Area Network (GAN) connects networks worldwide, offering seamless communication and data sharing on an international scale. It combines several WANs into a unified network infrastructure, enabling global connectivity.
	\end{NexMainBox}
	\begin{NexMainBox}[dark, crnA]
		\begin{NexListDark}
			\NexItemDark{\NexOption{Scope}: A GAN spans multiple continents, connecting regions and countries worldwide.}
			\NexItemDark{\NexOption{Applications}:
				\begin{NexListLight}
					\NexItemLight{Used by multinational corporations to synchronize operations across global offices.}
					\NexItemLight{Enables real-time communication and collaboration for global teams.}
					\NexItemLight{Supports international services like content delivery networks (CDNs) and global cloud platforms.}
				\end{NexListLight}
			}
			\NexItemDark{\NexOption{Technology}: Relies on satellite links, undersea cables, and high-speed WAN technologies to achieve global reach.}
		\end{NexListDark}
	\end{NexMainBox}
\end{NexMainBox}

\newpage
\subsection{Broadband Technologies}
\label{Broadband Technologies}
\begin{NexMainBox}[light, hdrA, sdwA, crnA, grwB, ssecA]
	\begin{NexMainBox}[dark, crnA]
		Broadband can be delivered through various technologies, each suited to specific user needs and geographic areas. Choosing the right technology depends on availability, cost, and required speed.
	\end{NexMainBox}
	\begin{NexMainBox}[dark, crnA]
		\begin{NexListDark}
			\NexItemDark{\NexLink{DSL (Digital Subscriber Line)}{DSL (Digital Subscriber Line)}: Provides internet through existing telephone lines, offering speeds up to 100 Mbps depending on distance and infrastructure.}
			\NexItemDark{\NexLink{Cable Broadband}{Cable Broadband}: Uses coaxial cables to deliver higher speeds than DSL, commonly ranging from 100 Mbps to 1 Gbps.}
			\NexItemDark{\NexLink{Fiber-Optic Broadband}{Fiber-Optic Broadband}: Utilizes light signals over fiber-optic cables for ultra-high speeds up to several Gbps, ensuring reliability and minimal latency.}
			\NexItemDark{\NexLink{Satellite Broadband}{Satellite Broadband}: Ideal for rural or remote areas, offering speeds up to 100 Mbps via satellite connections but with higher latency.}
			\NexItemDark{\NexLink{Mobile Broadband}{Mobile Broadband}: Delivers internet via cellular networks (3G, 4G, or 5G) with varying speeds and mobility advantages.}
		\end{NexListDark}
	\end{NexMainBox}
\end{NexMainBox}

\newpage
\subsubsection{DSL (Digital Subscriber Line)}
\label{DSL (Digital Subscriber Line)}
\begin{NexMainBox}[light, hdrA, sdwA, crnA, grwB, sssecA]
	\begin{NexMainBox}[dark, crnA]
		Digital Subscriber Line (DSL) is a broadband technology that delivers high-speed internet using existing copper telephone lines. It offers cost-effective connectivity and is widely available in urban and suburban areas.
	\end{NexMainBox}
	\begin{NexMainBox}[dark, crnA]
		\begin{NexListDark}
			\NexItemDark{\NexOption{Scope}: Limited by the distance between the user’s location and the telephone exchange, with speeds decreasing as the distance increases.}
			\NexItemDark{\NexOption{Applications}:
				\begin{NexListLight}
					\NexItemLight{Enables everyday browsing, streaming, and email communication.}
					\NexItemLight{Supports small business operations requiring stable internet.}
					\NexItemLight{Provides a solution for residential internet connectivity.}
				\end{NexListLight}
			}
			\NexItemDark{\NexOption{Technology}: Uses frequency division multiplexing (FDM) to separate voice calls and internet data over the same line.}
			\NexItemDark{\NexOption{Variants}:
				\begin{NexListLight}
					\NexItemLight{ADSL (Asymmetric DSL): Offers faster download speeds compared to upload speeds, ideal for web browsing and streaming.}
					\NexItemLight{VDSL (Very-high-bit-rate DSL): Provides higher speeds than ADSL, supporting activities like HD streaming and online gaming.}
					\NexItemLight{SDSL (Symmetric DSL): Ensures equal download and upload speeds, often used for business applications.}
				\end{NexListLight}
			}
		\end{NexListDark}
	\end{NexMainBox}
\end{NexMainBox}

\newpage
\subsubsection{Cable Broadband}
\label{Cable Broadband}
\begin{NexMainBox}[light, hdrA, sdwA, crnA, grwB, sssecA]
	\begin{NexMainBox}[dark, crnA]
		Cable Broadband is a popular high-speed internet technology that uses coaxial cables to transmit data. It offers faster speeds compared to DSL and is widely used in residential and commercial settings.
	\end{NexMainBox}
	\begin{NexMainBox}[dark, crnA]
		\begin{NexListDark}
			\NexItemDark{\NexOption{Scope}: Primarily serves homes and businesses in areas where cable television infrastructure is available.}
			\NexItemDark{\NexOption{Applications}:
				\begin{NexListLight}
					\NexItemLight{Enables streaming of high-definition videos and multimedia content.}
					\NexItemLight{Supports online gaming and real-time communication.}
					\NexItemLight{Facilitates faster downloads and uploads for personal and professional use.}
				\end{NexListLight}
			}
			\NexItemDark{\NexOption{Technology}: Utilizes DOCSIS (Data Over Cable Service Interface Specification) standards to deliver broadband services.}
			\NexItemDark{\NexOption{Advantages}:
				\begin{NexListLight}
					\NexItemLight{Provides higher speeds and reliability compared to DSL.}
					\NexItemLight{Supports multiple users and devices simultaneously.}
					\NexItemLight{Leverages existing cable TV infrastructure for widespread availability.}
				\end{NexListLight}
			}
			\NexItemDark{\NexOption{Limitations}:
				\begin{NexListLight}
					\NexItemLight{Shared bandwidth can cause speed fluctuations during peak usage.}
					\NexItemLight{Availability may be limited in rural or remote areas.}
					\NexItemLight{Requires a cable modem for connectivity.}
				\end{NexListLight}
			}
		\end{NexListDark}
	\end{NexMainBox}
\end{NexMainBox}

\newpage
\subsubsection{Fiber-Optic Broadband}
\label{Fiber-Optic Broadband}
\begin{NexMainBox}[light, hdrA, sdwA, crnA, grwB, sssecA]
	\begin{NexMainBox}[dark, crnA]
		Fiber-Optic Broadband is a cutting-edge internet technology that uses light signals transmitted through thin strands of glass or plastic (fiber-optic cables). It delivers ultra-fast and highly reliable internet speeds, making it ideal for modern digital demands.
	\end{NexMainBox}
	\begin{NexMainBox}[dark, crnA]
		\begin{NexListDark}
			\NexItemDark{\NexOption{Scope}: Designed to provide high-speed internet over large distances without signal degradation, suitable for residential, commercial, and industrial applications.}
			\NexItemDark{\NexOption{Applications}:
				\begin{NexListLight}
					\NexItemLight{Supports streaming ultra-high-definition (4K or 8K) videos seamlessly.}
					\NexItemLight{Enables real-time online gaming with minimal latency.}
					\NexItemLight{Facilitates large-scale data transfers for businesses and data centers.}
				\end{NexListLight}
			}
			\NexItemDark{\NexOption{Technology}: Utilizes pulses of light through fiber-optic cables to encode and transmit data at speeds up to several Gbps.}
			\NexItemDark{\NexOption{Advantages}:
				\begin{NexListLight}
					\NexItemLight{Provides unparalleled internet speeds and reliability.}
					\NexItemLight{Minimizes signal loss over long distances compared to copper cables.}
					\NexItemLight{Offers resistance to electromagnetic interference (EMI), ensuring stable connections.}
				\end{NexListLight}
			}
			\NexItemDark{\NexOption{Limitations}:
				\begin{NexListLight}
					\NexItemLight{Higher installation costs compared to other broadband technologies.}
					\NexItemLight{Limited availability in rural or underserved areas.}
					\NexItemLight{Requires professional installation due to the complexity of fiber-optic infrastructure.}
				\end{NexListLight}
			}
		\end{NexListDark}
	\end{NexMainBox}
\end{NexMainBox}

\newpage
\subsubsection{Satellite Broadband}
\label{Satellite Broadband}
\begin{NexMainBox}[light, hdrA, sdwA, crnA, grwB, sssecA]
	\begin{NexMainBox}[dark, crnA]
		Satellite Broadband provides internet access via communication satellites. It is especially beneficial for rural and remote areas where other broadband technologies are unavailable.
	\end{NexMainBox}
	\begin{NexMainBox}[dark, crnA]
		\begin{NexListDark}
			\NexItemDark{\NexOption{Scope}: Designed to cover areas where terrestrial broadband infrastructure is absent, including isolated regions and mobile platforms.}
			\NexItemDark{\NexOption{Applications}:
				\begin{NexListLight}
					\NexItemLight{Provides internet connectivity to remote homes, offices, and farms.}
					\NexItemLight{Supports communication for ships, aircraft, and other mobile platforms.}
					\NexItemLight{Enables emergency response services in disaster-affected regions.}
				\end{NexListLight}
			}
			\NexItemDark{\NexOption{Technology}: Relies on geostationary and low Earth orbit (LEO) satellites to transmit and receive signals. The user requires a satellite dish and modem for connectivity.}
			\NexItemDark{\NexOption{Advantages}:
				\begin{NexListLight}
					\NexItemLight{Provides internet access where other broadband technologies are unavailable.}
					\NexItemLight{Covers large geographical areas, including sparsely populated regions.}
					\NexItemLight{Supports continuous connectivity for mobile platforms.}
				\end{NexListLight}
			}
			\NexItemDark{\NexOption{Limitations}:
				\begin{NexListLight}
					\NexItemLight{Higher latency compared to terrestrial broadband due to signal travel distance.}
					\NexItemLight{Weather conditions can affect signal quality and reliability.}
					\NexItemLight{Typically more expensive than other broadband options due to equipment and service costs.}
				\end{NexListLight}
			}
		\end{NexListDark}
	\end{NexMainBox}
\end{NexMainBox}

\newpage
\subsubsection{Mobile Broadband}
\label{Mobile Broadband}
\begin{NexMainBox}[light, hdrA, sdwA, crnA, grwB, sssecA]
	\begin{NexMainBox}[dark, crnA]
		Mobile Broadband delivers internet connectivity through cellular networks such as 3G, 4G, and 5G. It enables on-the-go access for smartphones, tablets, and mobile hotspots.
	\end{NexMainBox}
	\begin{NexMainBox}[dark, crnA]
		\begin{NexListDark}
			\NexItemDark{\NexOption{Scope}: Ideal for personal or business use in urban, suburban, and even some rural areas with cellular coverage.}
			\NexItemDark{\NexOption{Applications}:
				\begin{NexListLight}
					\NexItemLight{Provides internet access for mobile devices like smartphones and tablets.}
					\NexItemLight{Facilitates remote work with portable hotspots.}
					\NexItemLight{Supports IoT devices and applications requiring constant connectivity.}
				\end{NexListLight}
			}
			\NexItemDark{\NexOption{Technology}: Operates on cellular networks such as LTE (4G) and 5G, offering varying speeds and coverage.}
			\NexItemDark{\NexOption{Advantages}:
				\begin{NexListLight}
					\NexItemLight{Provides mobility and flexibility, allowing users to stay connected anywhere with coverage.}
					\NexItemLight{Supports a wide range of devices beyond traditional computers.}
					\NexItemLight{Expands broadband access to underserved areas with cellular networks.}
				\end{NexListLight}
			}			\NexItemDark{\NexOption{Limitations}:
				\begin{NexListLight}
					\NexItemLight{Data caps and speed throttling may apply based on service plans.}
					\NexItemLight{Signal strength can vary depending on location and network congestion.}
					\NexItemLight{Latency may be higher compared to fiber or cable broadband.}
				\end{NexListLight}
			}
		\end{NexListDark}
	\end{NexMainBox}
\end{NexMainBox}


\section{About This Codebase}
This codebase aims to enhance the performance of the POSIX shell by offloading data manipulation tasks to powerful utilities like \texttt{awk}, \texttt{sed}, and \texttt{c90}.

\subsection{Key Features}
\begin{itemize}
    \item \textbf{Compatibility}: Adheres to POSIX draft 1003.2 (draft 11.3) standards.
    \item \textbf{Performance}: Utilizes \texttt{awk}, \texttt{sed}, and \texttt{c90} for backend processing.
    \item \textbf{Portability}: Scripts are designed to be portable and easily integrated into different environments without modification.
\end{itemize}
\newpage
\section{POSIX Nexus Overview}

The \texttt{run.sh} script initiates the POSIX Nexus daemon, which encompasses the following functionalities:

\begin{itemize}
    \item Database Management
    \item SSH Authentication
    \item Command Distribution across specified Nexuses
    \item Webserver functionalities akin to Cockpit, applicable to any system, not limited to \texttt{systemd}
    \item Comprehensive security measures ensuring information sharing among Nexuses and connected machines via multicast groups, authenticated through SSH keys
    \item Interactive menu for managing IP addresses of child machines, enabling secure connections through public keys
    \item Dual authentication mechanism ensuring both Nexus and group member validation, with an automatic ejection of any machine posing a security risk detected by multiple systems
\end{itemize}

The POSIX Nexus aims to create a robust network that prioritizes security without compromising usability. Routine tasks are handled via AWK, while advanced features are developed using the C90 standard library.

\section{Deployment Strategy}

Deployment will be executed incrementally, targeting supported systems and eventually expanding to all systems. This approach ensures thorough testing and stability, akin to Debian's release cycle methodology.

\section{Testing and Validation}

Testing will be conducted methodically, addressing each component individually to guarantee seamless integration and performance.

\section{Contribution Guidelines}
We welcome contributions from the community to help improve the project. Here are some guidelines to help you get started:

\subsection{Deployment Strategy}

Deployment will be executed incrementally, targeting supported systems and eventually expanding to all systems. This approach ensures thorough testing and stability, akin to Debian's release cycle methodology.

\subsection{Testing and Validation}

Testing will be conducted methodically, addressing each component individually to guarantee seamless integration and performance.

\subsection{Reporting Issues}
If you encounter any issues or have suggestions for improvements, please open an issue on our GitHub Issues page. Provide a clear and descriptive title, detailed description, steps to reproduce the issue, and any relevant logs or screenshots.

\subsection{Tool Restrictions}
To maintain compatibility and portability, it is prohibited to use any tools not explicitly required by POSIX. Ensure your scripts and contributions adhere strictly to POSIX standards.

\subsection{Testing}
Before submitting your pull request, please ensure your changes do not break existing functionality. Include tests for your changes if possible.

\subsection{Documentation}
If your contribution includes new features or changes to existing functionality, please update the documentation accordingly. Documentation files are located in the \texttt{docs} directory.

\subsection{Submitting Pull Requests}
We welcome pull requests for bug fixes, new features, and documentation improvements. Follow these steps to submit a pull request:
\begin{enumerate}
    \item \textbf{Fork the repository}: Click the "Fork" button at the top of the repository page to create a copy of the repository in your GitHub account.
    \item \textbf{Clone your fork}: Clone your forked repository to your local machine using the following command:
    \begin{lstlisting}
    git clone https://github.com/Canine-Table/posix-nexus.git
    \end{lstlisting}
\end{enumerate}

\subsection{Contact Us}
We are here to help and support the community. If you have any questions, concerns, or suggestions, please don't hesitate to reach out to us. Your feedback is invaluable in helping us improve and grow.

\begin{itemize}
    \item \textbf{Discord}: Join our  community on Discord for real-time discussions, support, and collaboration. [\href{https://discord.gg/GB9twwpCNM}{Discord Invite Link}]
    \item \textbf{GitHub Issues}: Report issues, request features, and track progress on our GitHub Issues page. [\href{https://github.com/Canine-Table/posix-nexus/issues}{GitHub Issues Link}]
    \item \textbf{Email}: For more detailed inquiries, feel free to email us at posix-nexus@duck.com.
\end{itemize}

\section{Submitting Pull Requests}
We welcome pull requests for bug fixes, new features, and documentation improvements. Follow these steps to submit a pull request:
\begin{enumerate}
    \item \textbf{Fork the repository}: Click the "Fork" button at the top of the repository page to create a copy of the repository in your GitHub account.
    \item \textbf{Clone your fork}: Clone your forked repository to your local machine using the following command:
    \begin{lstlisting}
    git clone https://github.com/Canine-Table/posix-nexus.git
    \end{lstlisting}
\end{enumerate}
\newpage
\end{document}
